\section{Conclusion}
\label{sec:conclusion}

This report treated the design of a world-action model as a space to be mapped
and built the instrument that maps it. The framework of \sref{sec:framework}
turns six independent decisions into interchangeable components behind one
training, deployment and evaluation interface, and it states each decision
formally, so a coupling architecture, a visibility pattern and a denoising
trajectory become objects with definite semantics. Two derivations carried the
design. Independent noise-level sampling across the two streams covers the whole
level square, which makes the deployment-time trajectory a variable a practitioner
selects after training. Cross-modal visibility determines whether an action-free
mixture component contributes a consistent video objective, which predicts that
the best visibility pattern changes once large-scale pretraining enters.

The study of \sref{sec:study} moved one axis at a time under one dataset, one
budget and one harness, and it selects a specific configuration. The action
stream becomes its own transformer that meets the video backbone through one
mixed attention at every layer. The backbone is the five-billion-parameter
text-and-image-to-video model, which balances accuracy against the cost of a
generation. The denoising target is a compact reconstructive latent, and a
frozen reducer brings a semantic latent to comparable quality by matching its
width, which identifies compactness as the property that matters. Visibility is
mutual once pretraining is available and one-directional before it. The corpus
co-trains human egocentric footage with robot trajectories in a single mixture,
which raises out-of-distribution success far more than it raises
in-distribution success. The two streams denoise together at deployment.

\sysa{} is that configuration scaled to a corpus of six thousand hours and
trained once, then fine-tuned per platform with nothing else changed. It reaches
the strongest reported world-action result on five of the nine simulated suites
in \sref{sec:sim}, and it operates three physical platforms including a
dexterous hand and three bimanual arm types in \sref{sec:real}. The deployment
stack of \sref{sec:cost} brings one generation into a range a thirty-hertz
controller absorbs, and the asynchronous executor removes the replanning stall
that a receding-horizon policy otherwise imposes on a physical loop.

Three directions stay open. Robustness to camera placement and to sensor noise
is the weakest axis of the current model, and the perturbation results of
\tref{tab:liberoplus} locate the gap precisely enough to target it with
viewpoint augmentation inside the pretraining mixture. Precision tasks remain
harder for the model than generalization tasks, which points at the temporal
resolution of the action chunk and at the granularity of the latent the video
stream denoises. The tri-system coupling reaches a higher success rate than the
selected one at a higher cost per generation, so a version of it that keeps the
understanding expert active only when the instruction demands semantic
disambiguation is worth building.

Every element of this report is released: the framework, the pretrained
checkpoint, the fine-tuned variant for each benchmark and each robot, and one
checkpoint per point of the design study. A reader who disagrees with a
conclusion can therefore reload the model that produced it and move the axis
again.
